% Generated draft: requires mathematical and source review.
\documentclass[11pt]{article}
\usepackage[margin=25mm]{geometry}
\usepackage{fontspec}
\setmainfont{Latin Modern Roman}
\usepackage{amsmath,amssymb,mathtools,unicode-math}
\setmathfont{Latin Modern Math}
\usepackage{enumitem,needspace,xurl,hyperref}
\setlength{\parindent}{0pt}
\setlength{\parskip}{5pt}
\setlength{\emergencystretch}{3em}
\newcommand{\N}{\mathbb{N}}
\newcommand{\Z}{\mathbb{Z}}
\newcommand{\Q}{\mathbb{Q}}
\newcommand{\R}{\mathbb{R}}
\newcommand{\C}{\mathbb{C}}
\newcommand{\F}{\mathbb{F}}
\newcommand{\A}{\mathbb{A}}
\newcommand{\PP}{\mathbb{P}}
\newcommand{\E}{\mathbb{E}}
\newcommand{\eps}{\varepsilon}
\newcommand{\dd}{\,\mathrm{d}}
\begin{document}
\section*{UnsolvedMath ID: 11000050}
\textbf{Source problem code:} catalogue entry
\subsection{Definitions and mathematical statement}
For \(g\geq3\), let \(H=H_1(\Sigma_g;\Z)\), let \(\mathcal I_g\) be the Torelli group, and let \(\mathcal K_g\) denote the Johnson kernel in the source's closed-surface convention. In the customary Johnson-homomorphism setting,
\[1\longrightarrow\mathcal K_g\longrightarrow\mathcal I_g
\xrightarrow{\pi}\bigl(\textstyle\bigwedge^3 H\bigr)/H\longrightarrow1,\]
where the copy of \(H\) uses the symplectic-form embedding \(h\mapsto h\wedge\omega\). This identifies the standard meaning suggested by the excerpt's reference to exact sequence (11); that sequence and its normalization are not reproduced.

Morita's question is whether
\[\forall g\geq3,\qquad H_1(\mathcal K_g;\Z)=\mathcal K_g/[\mathcal K_g,\mathcal K_g]\text{ is finitely generated as an abelian group}.\]
Equivalently, for each \(g\) there should be finitely many elements whose images generate the entire abelianization. This is integral finite generation, not just finite-dimensional rational homology, and is a question about \(\mathcal K_g\), not \(H_2(\mathcal I_g;\Q)\).

The source reports large abelian quotients constructed by Birman--Craggs--Johnson [BC, Jo1] and Morita [Mo2]. It proposes a negative-answer strategy suggested by [BF], Theorem 5.1, following [McM]: use the action of \(\mathcal K_g\) on \(H_1(Y;\Z)\) for a certain abelian cover \(Y\to\Sigma_g\), compatible with its deck-group module structure, to obtain
\[\rho:\mathcal K_g\to\operatorname{SL}_2(\Z[t,t^{-1}]),\]
then analyze its image and the associated Bruhat--Tits--Serre tree. Neither the cover, representation nor valuation/tree is specified here. The suggested ingredient [GMV] is reported to give free quotients of arbitrarily high rank for \(\operatorname{SL}_2(\Z[t])\) and \(\operatorname{SL}_2(K[s,t])\) for any finite field \(K\). That ambient-group report does not itself give such quotients of \(\rho(\mathcal K_g)\).

The historical paragraph says that for \(g\geq3\), \(\mathcal I_g\) is finitely generated and \(\mathcal K_g\) is not, and motivates studying \(\pi^{-1}(L)\) for subgroups \(L\leq(\bigwedge^3H)/H\). Those assertions about [BF] and \(\mathcal K_g\) must be verified against the exact source and subsequent corrections before being treated as present theorems. They are preserved as reports, not adopted as proof of an answer. Even non-finite-generation of a group alone would not decide finite generation of its abelianization. The pullbacks interpolate between \(\pi^{-1}(0)=\mathcal K_g\) and \(\pi^{-1}((\bigwedge^3H)/H)=\mathcal I_g\); the excerpt does not state a specific theorem for all intermediate subgroups.
\subsection{Short English statement}
Is the integral abelianization of the Johnson kernel finitely generated in every genus at least three? The source's tree-action strategy and historical finite-generation assertions are retained with explicit verification caveats.
\subsection{Sources}
\begin{itemize}
\item Local UnsolvedMath registry, permanent ID \(11000050\), exact source code \texttt{catalogue entry}.
\item Morita, attribution in Farb, \emph{Problems on Mapping Class Groups and Related Topics}, Question 5.2, PDF page 43; references to Theorem 5.1 and sequence (11): \url{https://www.math.uchicago.edu/~farb/papers/mcgbook.pdf}. Registry access 2026-07-29; pdf-font-structure extraction.
\item Internal references [BC, Jo1] (Birman--Craggs--Johnson), [Mo2] (Morita), [BF], [McM], and [GMV] (named in the source as Grunewald--Mennike--Vaserstein). Full bibliographic entries and precise theorem hypotheses are absent.
\item The research summary mentions unspecified 2023 results on \(H_2(\mathcal I_g;\Q)\) and abelianization work, without identifying a resolution of this exact question.
\end{itemize}
\end{document}
